\documentclass[10pt,a4paper]{article}
\usepackage[a4paper,margin=14mm,top=12mm,bottom=14mm]{geometry}
\usepackage{fontspec}
\setmainfont{Liberation Sans}
\usepackage{polyglossia}
\setdefaultlanguage{french}
\usepackage{amsmath}
\usepackage{xcolor}
\usepackage{graphicx}
\usepackage[most]{tcolorbox}
\usepackage{tabularx,booktabs,array,multirow}
\usepackage{enumitem}
\usepackage{fancyhdr}
\usepackage{lastpage}
\usepackage[hidelinks]{hyperref}
\definecolor{uimm}{HTML}{005677}
\definecolor{cyan}{HTML}{13A2DD}
\definecolor{fond}{HTML}{E6F2FA}
\definecolor{ok}{HTML}{2E7D32}
\definecolor{ko}{HTML}{B3261E}
\definecolor{violet}{HTML}{7B1FA2}
\setlength{\parindent}{0pt}
\setlength{\parskip}{3pt}
\setlist{nosep}
\setlist[itemize]{leftmargin=13pt,label=\textcolor{cyan}{\textbullet}}
\newcommand{\titre}[1]{\vspace{5pt}{\color{uimm}\large\bfseries #1}\par\vspace{-3pt}{\color{cyan}\rule{\linewidth}{1.2pt}}\par\vspace{0pt}}
\tcbset{boite/.style={colback=fond,colframe=uimm,boxrule=0.6pt,arc=2pt,left=6pt,right=6pt,top=4pt,bottom=4pt},
        defi/.style={colback=violet!6,colframe=violet,boxrule=0.8pt,arc=2pt,left=6pt,right=6pt,top=4pt,bottom=4pt,
                     title={\bfseries Défi},coltitle=white,colbacktitle=violet,fonttitle=\small}}
\newcommand{\bandeau}[3]{%
\begin{tcolorbox}[colback=uimm,colframe=uimm,arc=2pt,left=4pt,right=8pt,top=4pt,bottom=4pt]
\begin{minipage}[c]{18mm}\setlength{\fboxsep}{1.2pt}\fcolorbox{white}{white}{\includegraphics[height=16.5mm]{logo_uimm.jpg}}\end{minipage}\hspace{4mm}%
\begin{minipage}[c]{\dimexpr\linewidth-22mm\relax}\color{white}
{\small #1}\\[1pt]{\LARGE\bfseries #2}\\[2pt]{\large #3}\hfill{\small CQPM Technicien qualité}
\end{minipage}
\end{tcolorbox}}
\newcommand{\reponse}[1][1]{\par\vspace{2pt}{\color{gray!60}\foreach \i in {1,...,#1}{\rule{\linewidth}{0.3pt}\par\vspace{9pt}}}}
\usepackage{pgffor}
\pagestyle{fancy}\fancyhf{}
\renewcommand{\headrulewidth}{0pt}
\newcommand{\pied}[1]{\fancyfoot[L]{\footnotesize\color{uimm}#1}\fancyfoot[R]{\footnotesize\color{uimm}Page \thepage{} sur \pageref{LastPage}}}

\usepackage{tikz}
\usepackage{longtable}
\renewcommand{\bandeau}[3]{%
\begin{tcolorbox}[colback=uimm,colframe=uimm,arc=2pt,left=4pt,right=8pt,top=4pt,bottom=4pt]
\begin{minipage}[c]{18mm}\setlength{\fboxsep}{1.2pt}\fcolorbox{white}{white}{\includegraphics[height=16.5mm]{logo_uimm.jpg}}\end{minipage}\hspace{4mm}%
\begin{minipage}[c]{\dimexpr\linewidth-22mm\relax}\color{white}
{\small #1}\\[1pt]{\LARGE\bfseries #2}\\[2pt]{\large #3}
\end{minipage}
\end{tcolorbox}}
% en-tête d'une fiche : code, titre, question, bloc, durée, organisation, niveau
\newcommand{\entete}[7]{%
\begin{tcolorbox}[colback=uimm,colframe=uimm,arc=2pt,left=4pt,right=6pt,top=4pt,bottom=4pt]
\begin{minipage}[c]{18mm}\setlength{\fboxsep}{1.2pt}\fcolorbox{white}{white}{\includegraphics[height=16.5mm]{logo_uimm.jpg}}\end{minipage}\hspace{3mm}%
\begin{minipage}[c]{\dimexpr\linewidth-58mm\relax}\color{white}\raggedright
{\small Base SPC · #4}\\[1pt]{\Large\bfseries #1 · #2}\\[2pt]{\itshape #3}
\end{minipage}\hfill
\begin{minipage}[c]{33mm}\color{white}\small\raggedleft
Durée : \textbf{#5}\\ \textbf{#6}\\ Niveau #7 sur 3
\end{minipage}
\end{tcolorbox}}
\newcommand{\tagE}{{\setlength{\fboxsep}{1.5pt}\colorbox{cyan}{\color{white}\scriptsize\bfseries Ellistat}}}
\newcommand{\tagW}{{\setlength{\fboxsep}{1.5pt}\colorbox{uimm}{\color{white}\scriptsize\bfseries Page web}}}
\tcbset{pred/.style={colback=orange!6,colframe=orange!80!black,boxrule=0.6pt,arc=2pt,left=6pt,right=6pt,top=3pt,bottom=1pt,
                    title={\bfseries Je prédis, avant de commencer},coltitle=white,colbacktitle=orange!80!black,fonttitle=\small},
        retenir/.style={colback=fond,colframe=uimm,boxrule=0.6pt,arc=2pt,left=6pt,right=6pt,top=3pt,bottom=3pt,
                    title={\bfseries À retenir},coltitle=white,colbacktitle=uimm,fonttitle=\small},
        plus/.style={colback=violet!6,colframe=violet,boxrule=0.6pt,arc=2pt,left=6pt,right=6pt,top=3pt,bottom=3pt,
                    title={\bfseries Pour aller plus loin},coltitle=white,colbacktitle=violet,fonttitle=\small},
        act/.style={colback=uimm!10,colframe=uimm,boxrule=0.5pt,arc=2pt,left=5pt,right=5pt,top=2pt,bottom=2pt,before skip=7pt,after skip=3pt},
        verif/.style={colback=orange!6,colframe=orange!80!black,boxrule=0.6pt,arc=2pt,left=6pt,right=6pt,top=3pt,bottom=3pt,
                    title={\bfseries À vérifier sur poste avant la première séance},coltitle=white,colbacktitle=orange!80!black,fonttitle=\small}}
\newcommand{\fichier}[1]{\texttt{\small #1}}
\newcommand{\menu}[1]{\textsf{\textbf{#1}}}
\newcommand{\champ}[1]{\rule{#1}{0.3pt}}

\pied{Base SPC · fiche A2 · Ma première carte p, à la main}
\begin{document}
\thispagestyle{fancy}
\entete{A2}{Ma première carte p, à la main}{Le taux de pièces défectueuses est-il stable ?}{Cartes aux attributs}{35 min}{en binôme}{1}
\begin{tcolorbox}[boite]\textbf{Le but.} Relever 20 lots, calculer la moyenne et les limites, tracer la carte à la main.\par\textbf{Il vous faut.} La page « Tirer un lot » (scénario donné par le formateur) · Une calculatrice, un crayon, une règle\end{tcolorbox}
\begin{tcolorbox}[pred]Dans des lots de 50 pièces, on trouve en moyenne 2 défectueuses. Un lot avec 5 défectueuses est-il anormal ?\reponse\end{tcolorbox}
\titre{À faire}
\begin{enumerate}[label=\textcolor{cyan}{\textbf{\arabic*.}},leftmargin=15pt,itemsep=1.5pt]
\item \tagW\ Page « Tirer un lot » : on compte des pièces défectueuses, taille du lot 50, scénario et numéro de poste donnés par le formateur.
\item Tirer 20 lots, un par un. Noter chaque résultat dans le tableau et placer le point sur la carte.
\item Calculer $\bar p$ = total des défectueuses / total des pièces contrôlées.
\item Calculer les limites : $\bar p \pm 3\sqrt{\bar p\,(1-\bar p)/n}$ avec $n = 50$. Une limite négative devient 0.
\item Tracer la ligne centrale et les limites sur la carte (en nombre de pièces : multiplier par 50).
\end{enumerate}
\vspace{3pt}{\renewcommand{\arraystretch}{1.35}\small\begin{tabularx}{\linewidth}{|l|*{10}{>{\centering\arraybackslash}X|}}\hline
\textbf{Lot} & 1 & 2 & 3 & 4 & 5 & 6 & 7 & 8 & 9 & 10\\\hline
\textbf{Défectueuses} &  &  &  &  &  &  &  &  &  & \\\hline\end{tabularx}}\par\vspace{2pt}
{\renewcommand{\arraystretch}{1.35}\small\begin{tabularx}{\linewidth}{|l|*{10}{>{\centering\arraybackslash}X|}}\hline
\textbf{Lot} & 11 & 12 & 13 & 14 & 15 & 16 & 17 & 18 & 19 & 20\\\hline
\textbf{Défectueuses} &  &  &  &  &  &  &  &  &  & \\\hline\end{tabularx}}\par\vspace{2pt}
\begin{center}\begin{tikzpicture}[x=0.780cm,y=0.31cm]\draw[gray!45,very thin] (0,0) grid (20,12);\draw[thick] (0,0)--(20,0);\draw[thick] (0,0)--(0,12);\node[left,font=\scriptsize] at (0,0) {0};\node[left,font=\scriptsize] at (0,2) {2};\node[left,font=\scriptsize] at (0,4) {4};\node[left,font=\scriptsize] at (0,6) {6};\node[left,font=\scriptsize] at (0,8) {8};\node[left,font=\scriptsize] at (0,10) {10};\node[left,font=\scriptsize] at (0,12) {12};\node[below,font=\scriptsize] at (1,0) {1};\node[below,font=\scriptsize] at (2,0) {2};\node[below,font=\scriptsize] at (3,0) {3};\node[below,font=\scriptsize] at (4,0) {4};\node[below,font=\scriptsize] at (5,0) {5};\node[below,font=\scriptsize] at (6,0) {6};\node[below,font=\scriptsize] at (7,0) {7};\node[below,font=\scriptsize] at (8,0) {8};\node[below,font=\scriptsize] at (9,0) {9};\node[below,font=\scriptsize] at (10,0) {10};\node[below,font=\scriptsize] at (11,0) {11};\node[below,font=\scriptsize] at (12,0) {12};\node[below,font=\scriptsize] at (13,0) {13};\node[below,font=\scriptsize] at (14,0) {14};\node[below,font=\scriptsize] at (15,0) {15};\node[below,font=\scriptsize] at (16,0) {16};\node[below,font=\scriptsize] at (17,0) {17};\node[below,font=\scriptsize] at (18,0) {18};\node[below,font=\scriptsize] at (19,0) {19};\node[below,font=\scriptsize] at (20,0) {20};\node[right,font=\scriptsize] at (20,0) {0 \%};\node[right,font=\scriptsize] at (20,2) {4 \%};\node[right,font=\scriptsize] at (20,4) {8 \%};\node[right,font=\scriptsize] at (20,6) {12 \%};\node[right,font=\scriptsize] at (20,8) {16 \%};\node[right,font=\scriptsize] at (20,10) {20 \%};\node[right,font=\scriptsize] at (20,12) {24 \%};\node[rotate=90,font=\scriptsize] at (-1.41,6.0) {pièces défectueuses};\node[below,font=\scriptsize] at (10.0,-1.3) {lot};\end{tikzpicture}\end{center}
\newpage
\titre{J'observe, je décide}
\begin{enumerate}[label=\textcolor{cyan}{\textbf{Q\arabic*.}},leftmargin=22pt]
\item $\bar p$ = \champ{16mm} \% \quad Limite supérieure = \champ{16mm} \% soit \champ{12mm} pièces \quad Limite inférieure = \champ{16mm}\par\vspace{5pt}
\item Un point sort-il des limites ? Le procédé est-il stable ?\reponse[1]
\item Reprenez la prédiction : 5 défectueuses dans un lot de 50, est-ce anormal ?\reponse[1]
\item Ouvrir « Afficher la carte calculée par la page ». Vos limites sont-elles les mêmes ?\reponse[1]
\end{enumerate}
\begin{tcolorbox}[retenir]Les limites disent ce que le hasard produit à lui seul. Tant que les points restent dedans, on ne touche à rien.\end{tcolorbox}
\begin{tcolorbox}[plus]Sans écran : un lot = 10 dés lancés, une pièce défectueuse = un 6. Relever 20 lancers et tracer la carte avec $n = 10$.\end{tcolorbox}
\vfill{\footnotesize\color{uimm}\textbf{Ressources de la base :} Tirer un lot · Jeu de secours.}
\end{document}